% TOP_PROBLEM: {"id": 2772, "title": "Kirby Problem 2.24", "category_id": 7, "category": {"id": 7, "name": "topology", "display_name": "Topology", "description": "Properties preserved under continuous deformations.", "slug": "topology", "order_index": 7, "created_at": "2026-07-31T15:26:25.671Z"}, "status": "open", "problem_number": "KP-2.24", "set_id": 11, "statusQualification": "Distinct fiberings are counted by different fiber subgroups of the fixed total-space fundamental group, as in the cited surface-bundle literature."}
% FIRST_POSTED: 2026-10-01
% ENTRY_KIND: statement_only
% UnsolvedMath ID 2772; source catalogue code KP-2.24
% !TeX program = lualatex
% Definition and sources only; this file is not an LLM solution attempt.
\documentclass[11pt,a4paper]{article}
\usepackage[margin=25mm]{geometry}
\usepackage{fontspec}
\setmainfont{lmroman10-regular.otf}[BoldFont=lmroman10-bold.otf,
 ItalicFont=lmroman10-italic.otf,BoldItalicFont=lmroman10-bolditalic.otf]
\usepackage{amsmath,amssymb,mathtools}
\usepackage{unicode-math}
\setmathfont{latinmodern-math.otf}
\AtBeginDocument{\let\setminus\smallsetminus}
\usepackage{xurl,hyperref}
\hypersetup{colorlinks=true,urlcolor=blue}
\setcounter{secnumdepth}{0}
\setlength{\parindent}{0pt}
\setlength{\parskip}{5pt}
\setlength{\emergencystretch}{3em}
\begin{document}
\section*{Kirby Problem 2.24}\label{2772}
{\small UnsolvedMath ID 2772 · KP-2.24 · Topology · Kirby's Problems in Low-Dimensional Topology}
\subsection{Definitions and mathematical statement}
A complex surface is a connected complex manifold of complex dimension two.
Consider a compact such surface $X$. A surface-bundle structure on its underlying
smooth four-manifold is a smooth locally trivial map $p:X\to B$ whose base $B$
and fiber $F$ are closed connected oriented surfaces of genus at least two.
Local triviality means that over each sufficiently small open set $U\subset B$
there is a diffeomorphism $p^{-1}(U)\cong U\times F$ commuting with projection.

Use the fiber-subgroup convention for distinct fiberings in the cited
surface-bundle literature. A fibration determines the normal subgroup
$K_p=\ker(p_*:\pi_1(X)\to\pi_1(B))$, identified with the fiber group.
Two fiberings are equivalent in this convention when their fiber subgroups
are equal inside $\pi_1(X)$. Normality removes basepoint-conjugacy ambiguity.
Reparametrizing the base does not change the subgroup; arbitrary
automorphisms of the total space are not used to identify different subgroups.

Does there exist a compact complex surface with at least three surface-bundle
structures having pairwise different fiber subgroups? The same complex surface
must support all the structures; constructing unrelated bundles with diffeomorphic
fibers is insufficient. The genus restrictions are those in the source discussion:
without them, torus constructions change the problem substantially. The question
concerns bundles on a complex surface, without adding a requirement that their
local trivializations be holomorphic.
\subsection{Short English statement}
Can one compact complex surface fiber over surfaces in at least three different ways, distinguished by their fiber subgroups, with both base and fiber of genus at least two?
\subsection{Sources}
\begin{itemize}
\item[S1] ULAM.AI and the UnsolvedMath Contributors, supplied problems.json, record 2772 (KP-2.24), Kirby's Problems in Low-Dimensional Topology. Catalogue identity and source transcription; CC BY 4.0.
\url{https://huggingface.co/datasets/ulamai/UnsolvedMath}.
\item[S2] K3: A New Problem List in Low-Dimensional Topology, Problem 2.24. Expanded formulation of the supplied catalogue statement.
\url{https://math.berkeley.edu/sites/default/files/surv-295-ruberman-watermarked-author-pdf.pdf}.
\item[S3] N. Salter, Surface bundles over surfaces with arbitrarily many fiberings, Section 2 and Proposition 2.2. Fiber-subgroup equivalence convention.
\url{https://nsalter.science.nd.edu/research/multifiberedsbs.pdf}.
\end{itemize}
\end{document}
